\documentclass[11pt]{article}

\usepackage[T1]{fontenc}
\usepackage{amsmath,amssymb,amsthm}
\usepackage[margin=1in]{geometry}

\newtheorem{theorem}{Theorem}

\newcommand{\N}{\mathbb{N}}

\title{Why the Ordinary-Dice Game Is Never Fair}
\author{}
\date{}

\begin{document}

\maketitle

Player A rolls $a$ ordinary fair six-sided dice, and Player B rolls $b$.
A die scores one point exactly when it shows a six, and otherwise scores no
points.  The player with the higher score wins, with ties awarded to B.

\begin{theorem}
If $a,b\in\N$ are positive, then the probability that B wins is not
$\tfrac12$.
\end{theorem}

\begin{proof}
For $0\leq k\leq n$, let
\[
  c(n,k)=\binom{n}{k}5^{n-k}
\]
be the number of outcomes in which exactly $k$ of $n$ dice show a six.  The
binomial coefficient chooses the successful dice, while each of the remaining
$n-k$ dice has five possible faces.

Let $W(a,b)$ denote the number of elementary outcomes in which B wins.  If A
scores $i$, then B must score some $j\geq i$, so
\[
  W(a,b)
    =\sum_{i=0}^{a}\ \sum_{j=i}^{b} c(a,i)c(b,j).
\]
(An inner sum is empty when $i>b$.)  Since all $6^{a+b}$ joint die outcomes
are equally likely,
\[
  \Pr(B\text{ wins})=\frac{W(a,b)}{6^{a+b}}.
\]

We now compute the winning count modulo $5$.  From the definition of $c$,
\[
  c(n,k)\equiv
  \begin{cases}
    1 & \text{if }k=n,\\
    0 & \text{if }k\neq n
  \end{cases}
  \pmod 5.
\]
Indeed, if $k<n$, the positive power $5^{n-k}$ supplies a factor of $5$; if
$k=n$, then $c(n,n)=1$.  Values with $k>n$ contribute nothing because the
corresponding binomial coefficient is zero.

Thus every summand defining $W(a,b)$ vanishes modulo $5$ except possibly the
one with $i=a$ and $j=b$.  That term occurs precisely when $a\leq b$, and
hence
\[
  W(a,b)\equiv
  \begin{cases}
    1 & \text{if }a\leq b,\\
    0 & \text{if }a>b
  \end{cases}
  \pmod 5.                                      \tag{1}
\]

Suppose, for a contradiction, that the game is fair.  Then
\[
  \frac{W(a,b)}{6^{a+b}}=\frac12,
  \qquad\text{so}\qquad
  W(a,b)=\frac{6^{a+b}}2.
\]
Because $a+b>0$, we can take one factor of $6$ out of the power and obtain
\[
  W(a,b)=3\cdot 6^{a+b-1}.
\]
Reducing this identity modulo $5$ gives
\[
  W(a,b)\equiv 3\cdot 1^{a+b-1}\equiv 3\pmod 5,
\]
since $6\equiv1\pmod5$.  This contradicts (1), which says that the same count
is congruent to either $0$ or $1$ modulo $5$.  Therefore the probability that
B wins cannot be $\tfrac12$.
\end{proof}

\end{document}
